\documentclass[12pt]{article}

% Idioma e codificação
\usepackage[brazil]{babel}
\usepackage[utf8]{inputenc}
\usepackage[T1]{fontenc}

% Matemática e listas
\usepackage{amsmath, amssymb}
\usepackage{enumitem}

% Layout e cores
\usepackage{geometry}
\usepackage{graphicx}
\usepackage{xcolor}

% Cores institucionais
\definecolor{maincolor}{RGB}{0,100,50}
\definecolor{forestgreen}{RGB}{0,100,50}

% Cabeçalho e rodapé
\usepackage{fancyhdr}

% Blocos
\usepackage[most]{tcolorbox}
\newtcolorbox{block}[1]{%
  colback=green!5,
  colframe=maincolor,
  coltitle=white,
  title=#1,
  fonttitle=\bfseries,
  boxrule=1.5pt,
  arc=3mm,
  left=6pt,
  right=6pt,
  top=6pt,
  bottom=6pt
}

% Margens
\geometry{margin=2.5cm}

% Fancyhdr
\pagestyle{fancy}
\fancyhf{}

\renewcommand{\headrulewidth}{2pt}
\renewcommand{\footrulewidth}{2pt}

\renewcommand{\headrule}{\hbox to\headwidth{%
  \color{forestgreen}\leaders\hrule height \headrulewidth\hfill}}

\renewcommand{\footrule}{\hbox to\headwidth{%
  \color{forestgreen}\leaders\hrule height \footrulewidth\hfill}}

\setlength{\headheight}{52pt}
\addtolength{\topmargin}{-20pt}

\fancyhead[L]{%
  \textbf{UniRV} \\ \small Universidade de Rio Verde
}

\fancyhead[R]{%
  \raggedleft
  \textcolor{forestgreen}{%
    \small
    \textbf{Lista de Exercícios: Métricas Google SRE}\\
    Me.~Sergio Souza Novak\\
    \textit{Engenharia da Qualidade e Confiabilidade}
    }
}

\fancyfoot[C]{\textcolor{forestgreen}{\thepage}}

\newtcolorbox{resposta}[1]{%
  colback=white,
  colframe=forestgreen!30,
  arc=1mm,
  height=#1,
  width=\textwidth,
  boxrule=0.8pt,
  breakable
}

\begin{document}

\begin{center}
  {\Large \textbf{Lista de Exercícios: Métricas Google SRE}}\\[0.2cm]
  \textit{Engenharia de Software / Engenharia da Qualidade}
\end{center}

\vspace{0.4cm}

\section*{Parte 6 — Métricas de Requisição e Confiabilidade em APIs}

\textbf{16.}
\begin{resposta}{3.5cm}
Disponibilidade baseada em tempo mede quanto tempo o sistema ficou operacional em relação ao tempo total (ex: 99,9\% ao mês).  
Disponibilidade baseada em requisições mede quantas requisições foram atendidas com sucesso em relação ao total recebido.  
A primeira foca em janela temporal; a segunda foca na experiência real do usuário.
\end{resposta}

\textbf{17.}
\begin{resposta}{3.5cm}
Erros 4xx indicam problema na requisição do cliente (ex: requisição inválida ou não autorizada).  
Não representam falha do servidor.  
Por isso, normalmente apenas erros 5xx são considerados falhas do sistema para cálculo de disponibilidade.
\end{resposta}

\textbf{18.}
\begin{resposta}{5cm}
Total = 1.200.000  

a) Apenas 2xx como sucesso:  
Disponibilidade = 1.199.400 / 1.200.000 = 0,9995 = 99,95\%.

b) Considerando apenas 5xx como falha:  
Falhas reais = 200  
Disponibilidade = (1.200.000 - 200) / 1.200.000  
= 1.199.800 / 1.200.000  
= 0,999833 = 99,9833\%.
\end{resposta}

\textbf{19.}
\begin{resposta}{5cm}
Meta = 99,9\%  

a) Erro permitido = 0,1\%  
0,001 × 2.500.000 = 2.500 falhas permitidas.

b) Como ocorreram 1.200 falhas 5xx,  
1.200 < 2.500  
Logo, a meta foi cumprida.
\end{resposta}

\textbf{20.}
\begin{resposta}{6cm}
Total = 6  

a) Apenas 2xx como sucesso:  
3 / 6 = 0,5 = 50\%.

b) Apenas 5xx como falha real:  
Falha real = 1  
Disponibilidade = 5 / 6 = 83,33\%.

c) Em produção, considera-se apenas 5xx como falha do sistema, pois 4xx representam erro do cliente.
\end{resposta}

\textbf{21.}
\begin{resposta}{3.5cm}
POFOD = Número de falhas / Número de solicitações  

POFOD = 350 / 700.000  
= 0,0005  

Logo, probabilidade de falha sob demanda = 0,05\%.
\end{resposta}

\textbf{22.}
\begin{resposta}{5cm}
a) MTTF = 1 / ROCOF  
MTTF = 1 / 0,002 = 500 horas.

b) MTTR = 30 minutos = 0,5 hora  

Disponibilidade = MTTF / (MTTF + MTTR)  
= 500 / (500 + 0,5)  
= 500 / 500,5  
= 0,999 = 99,9\%.
\end{resposta}

\textbf{23.}
\begin{resposta}{4cm}
Error Budget é a quantidade de falhas permitidas dentro da meta de disponibilidade.  
Para 99,9\%, o sistema pode falhar 0,1\% do tempo ou das requisições.  
Quanto maior a disponibilidade (99,99\% ou 99,999\%), menor o orçamento de erro e maior o rigor operacional.
\end{resposta}

\textbf{24.}
\begin{resposta}{4cm}
Em sistemas industriais, Confiabilidade é mais crítica.  
Uma falha pode causar acidentes ou prejuízos graves.  
É preferível o sistema parar (baixa disponibilidade) do que operar incorretamente.
\end{resposta}

\end{document}